\chapter{Limitations and Perspectives}
\label{chap:perspectives}

Although Chapter~\ref{chap:results} establishes controlled evidence, transfer,
repeatability and recovery remain unresolved. The discussion below links each
limitation to its engineering consequence and to the evaluation needed before
broader claims can be made.

\section{Measured limits}

\paragraph{Simulation and operation.}
The six-robot check was one cold start, not a simultaneous-motion or
repeatability study. Room~315 uses a kinematic shuttle backend that omits torque,
slip, braking, play and payload inertia. Positive and fault cases were each run
once; long
outages, profile switching, longer wait-for cycles, temporal non-progress and
device acknowledgement correlated to a specific command remain untested.

\paragraph{Perception and data.}
The visual results concern the known fleet, fixed simulated cameras and
generated scene families. Scene IDs, generation-payload fingerprints, image
bytes and complete label rows are disjoint, but a post-hoc audit found repeated
target-state, trajectory and geometry support across development partitions.
This is not sample reuse, yet it means the evidence is not independent of the
underlying synthetic distribution. The Final Test therefore says nothing about
new hardware, lighting, textures or installation conditions. Identity presence
and rail side also come from configured registries rather than learned outputs.

The segment score is Validation-calibrated, the loaded-state score is not a
calibrated probability, and the labelled 0.05/0.12 tolerances are offline
measures rather than runtime confidence. Their relationship to the separate
online slot-radius rule is fixed in \cref{tab:visual-threshold-map}. The
unexplained 52/55 shadow acceptance is an additional observation-level limit.

\paragraph{Task and interface scope.}
The twelve command-producing cases cover selected valid scenes only; neither
the Final Test nor the observation studies execute a task. English evaluation
uses ten selected requests, so open-vocabulary and multilingual behaviour are
unknown. Finally, CLI risk labels are heuristics and the subscriber assumes a
trusted ROS graph; multi-user use would require authenticated sources or graph
isolation.

\section{Priority next steps}

\paragraph{Evaluation order and stopping rule.}
The completed Final Test remains a fixed reference and must not become a tuning
set. I would first create the domain-shift data and repeat the existing
observation and campaign protocols unchanged. A failure would open a new
development cycle evaluated on Validation and the post-selection regression
set; it would not justify modifying the completed Test result. Physical
deployment is governed by the separate trial gate defined below.

\begin{table}[H]
\centering
\small
\renewcommand{\arraystretch}{1.08}
\caption{Priority experiments for extending the evidence.}
\label{tab:remaining-plan}
\begin{tabularx}{\textwidth}{@{}p{0.10\textwidth}p{0.36\textwidth}X@{}}
\toprule
\textbf{Priority} & \textbf{Engineering task} & \textbf{Question answered}\\
\midrule
P0 & Create a separate visual-shift set with varied cameras, lighting, textures
and sensor noise & Does perception transfer beyond the completed synthetic
distribution without tuning on the Final Test?\\
\tablerowrule
P0 & Repeat the seven-scene observation review across declared seeds & How
variable are acceptance and localisation in sparse, dense and blocker scenes,
including the current 52/55 shadow result?\\
\tablerowrule
P0 & Inject timed camera, presence, planner and node-restart faults in the live
ROS graph & Can the runtime recover while preserving fail-closed behaviour?\\
\tablerowrule
P0 & Repeat the positive and fault campaigns across controlled seeds & How
variable are completion, abstention, recovery and controller stop?\\
\tablerowrule
P1 & Correlate switch/stopper acknowledgements with commands and extend
reservation/deadlock cases & Do devices and multiple shuttles remain consistent
under delayed feedback and longer conflicts?\\
\tablerowrule
P1 & Mirror the releases and perform a clean-machine restore and replay & Can
MFJA recover the exact source, data and evidence independently?\\
\bottomrule
\end{tabularx}
\end{table}

\section{Continued use at MFJA}

The handover remains suitable for Gazebo teaching and research: students can
launch selected robots, inspect interfaces, alter rail scenes, generate PDDL
and trace a supervised task. Alternative planners, perception checkpoints or
grounding methods can reuse the fixed cases without changing controller or
arrival contracts. New model selection must use Validation; the small
post-selection regression set remains a development check, and new visual
conditions require a separate domain-shift evaluation rather than reuse of the
completed Final Test. Here, domain shift means a change in the visual-data
distribution, not a change to the PDDL action domain or ROS Domain ID.

\paragraph{Gate before a physical trial.}
Simulation qualification would be a prerequisite, not deployment authority. A
future physical trial would also require an independent machinery-risk review,
an explicit mapping to the platform's certified emergency functions, an
authenticated command owner, and a low-speed commissioning protocol with
predefined abort criteria. None of those controls is supplied by the current
Gazebo supervisor, so a successful simulation replay alone must not be used to
authorise hardware motion.

The completed Final Test should remain unchanged. Future experiments should
examine new visual conditions, repeated runs, live faults and concurrent shuttle
operation.
